\begin{prob}
    % Write the problem below. Precede any paths with \parentdir. For example:
    %   \includegraphics{\parentdir/public/fig/picture.png}
    % Any files used in the solution should be placed in the `private` 
    % subdirectory. Any other files can go in the `public` subdirectory.

    Determine whether each piece of code is correct or
    incorrect by thinking inductively. If it is correct, say so (you do not
    need to show work). But if it is incorrect, say what will go wrong (e.g.,
    infinite recursion, incorrect answer), give an example of an input that
    will demonstrate the error, and briefly explain \emph{why} the code does not
    work correctly.

    Hint: use the tips for analyzing recursive arguments from lecture.


    \begin{subprobset}
        \begin{subprob}
            In this part, adopt the convention that the sum of the numbers in an empty
            list is zero.

            \begin{minted}[autogobble]{python}
                import math
                def summation(numbers):
                    """Given a list, returns the sum of the numbers in the list."""
                    n = len(numbers)
                    if n == 0:
                        return 0
                    middle = math.floor(n / 2)
                    return summation(numbers[:middle]) + summation(numbers[middle:])
            \end{minted}

            \begin{soln}
                This is missing a base case for when the array is of size 1. If
                an input of size one is given to this function, it will make a
                recursive call on an empty array and an array of size one. This
                recursive call will make the same recursive calls, \emph{ad
                infinitum}. In other words, this will recurse infinitely.
            \end{soln}
        \end{subprob}

        \begin{subprob}
            In this part, you may assume that the input list is not empty.

            \begin{minted}[autogobble]{python}
                def product(numbers):
                    """Returns the product of all elements in numbers."""
                    if len(numbers) == 1:
                        return numbers[0]
                    return numbers[0] * product(numbers[1:])
            \end{minted}

            \begin{soln}
                This is correct.
            \end{soln}
        \end{subprob}

        \begin{subprob}
            You may assume that \python{start} and \python{stop} are valid indices.
            \begin{minted}[autogobble]{python}
                def inplace_reverse(l, start, stop):
                    """Reverses l[start:stop], in place."""
                    if stop - start <= 1:
                        return
                    l[start], l[stop-1] = l[stop-1], l[start] # swap first and last
                    inplace_reverse(l, start, stop-1)
            \end{minted}

            \begin{soln}
                The arguments to the recursive call are wrong: they should sort the middle part
                of the list (excluding the first and last element), but they include the first
                element. This causes incorrect output. For instance, \python{inplace_reverse([4,5,2,3], 0, 4)} results in
                \python{[3,5,2,4]}.
            \end{soln}
        \end{subprob}


        \begin{subprob}
            You may assume that \python{arr} is non-empty.

                \begin{minted}[autogobble]{python}
                    def find_mode(arr):
                    """Returns the mode of arr along with its frequency."""
                        if len(arr) == 1:
                            return arr[0], 1
                        middle = math.floor(len(arr) / 2)
                        left_mode, left_freq = find_mode(arr[:middle])
                        right_mode, right_freq = find_mode(arr[:middle])
                        if left_freq > right_freq:
                            return left_mode, left_freq
                        else:
                            return right_mode, right_freq
                \end{minted}

                \begin{soln}
                    This will return the incorrect answer. Assuming that recursive calls
                    work correctly, the algorithm may still not produce the right result.
                    For instance, suppose our input is:
                    \python{[1, 1, 1, 2, 2, 3, 3, 3, 2, 2]}.
                    The mode is clearly \python{2}, but the recursive calls will find
                    \python{1} and \python{3} for the mode in each half.

                    In short, the overall mode is not necessarily the mode of the left
                    half or the right half.
                \end{soln}

        \end{subprob}



    \end{subprobset}

\end{prob}
